\clearpage
\section{Технологическая часть}

\subsection{Средства реализации и интерфейс}

Алгоритмы реализованы на C++17. Для графа используется симметричная матрица пар «время, стоимость». Суммы хранятся в 64-разрядном целочисленном типе. При 200 портах и метках до $10^9$ путь содержит до 199 рёбер, поэтому его суммарное время и стоимость представимы в этом типе.

Входной файл содержит $n$, лимит $B$, $n$ имён портов и пары меток рёбер верхнего треугольника матрицы в порядке строк. Пара \texttt{0 0} обозначает отсутствующее ребро. Программа принимает режимы \texttt{exact}, \texttt{ant} и \texttt{both}. Для полного перебора установлен предел в 11 портов. В результате указываются последовательность портов, время маршрута, его стоимость и длительность поиска. Статус \texttt{infeasible} означает доказанное отсутствие решения, а \texttt{not\_found} --- исчерпание попыток эвристики.

\subsection{Реализации алгоритмов}

Листинг~\ref{lst:exact} содержит перебор перестановок индексов портов.

\lstinputlisting[firstline=93,lastline=117,firstnumber=93,
  basicstyle=\ttfamily\footnotesize\setstretch{1},
  caption={Реализация полного перебора},label={lst:exact}]{src/algorithms.cpp}

\clearpage
Для случайного выбора используется генератор \texttt{std::mt19937\_64}. Внутри построения маршрута список кандидатов и веса рассчитываются заново для текущего порта. Вещественные параметры ограничены значениями $\alpha,\beta\in[0,5]$ и $\rho\in[0,1]$. При допустимых метках графа и концентрациях феромона исходные веса представимы положительными числами в типе \texttt{double}. В листинге~\ref{lst:walk} показано, как веса преобразуются в вероятности по формуле~\eqref{eq:probability}.

\lstinputlisting[firstline=50,lastline=65,firstnumber=50,
  basicstyle=\ttfamily\footnotesize\setstretch{1},
  caption={Вычисление весов и вероятностей перехода},label={lst:walk}]{src/algorithms.cpp}

В листинге~\ref{lst:deposit} показано, как добавки к феромону накапливаются симметрично для обоих направлений перехода.
\lstinputlisting[firstline=79,lastline=90,firstnumber=79,
  basicstyle=\ttfamily\footnotesize\setstretch{1},
  caption={Накопление добавок феромона},label={lst:deposit}]{src/algorithms.cpp}

\clearpage
Основной цикл муравьиного алгоритма приведён в листинге~\ref{lst:ant}. Лучший результат сохраняется между итерациями.
\lstinputlisting[firstline=119,lastline=153,firstnumber=119,
  basicstyle=\ttfamily\footnotesize\setstretch{1},
  caption={Основной цикл муравьиного алгоритма},label={lst:ant}]{src/algorithms.cpp}

\clearpage
\subsection{Функциональное тестирование}

Для проверки полного перебора отдельная программа на Python генерировала все перестановки и вычисляла минимум по паре $(L,C)$. Проверялись также число портов в маршруте, отсутствие повторов, наличие каждого ребра и соблюдение бюджета. Для эвристики дополнительно проверена повторяемость результатов при одинаковом начальном состоянии генератора. В таблице~\ref{tab:cases} приведены примеры ручных тестов.

\begin{table}[H]\centering
\caption{Примеры проверок}
\label{tab:cases}
\begin{tabularx}{\linewidth}{|>{\raggedright\arraybackslash}p{0.57\linewidth}|>{\raggedright\arraybackslash}X|}
\hline
\firsttabhead{Входные данные} & \tabhead{Ожидаемый результат} \\
\hline
Один порт, $B=0$ & $L=0$, $C=0$ \\
Одно ребро $(4,9)$, $B=9$ & $L=4$, $C=9$ \\
Треугольник с рёбрами $(1,8)$, $(6,2)$, $(2,2)$, $B=4$ & $L=8$, $C=4$ \\
Тот же граф, $B=12$ & $L=3$, $C=10$ \\
Цепь с рёбрами $(5,3)$ и $(7,4)$, $B=100$ & $L=12$, $C=7$ \\
Граф с изолированным портом & Допустимых маршрутов нет \\
Указанный треугольник, $B=3$ & Допустимых маршрутов нет \\
\hline
\end{tabularx}
\end{table}

Все 20 проверок с независимым эталоном и 6 проверок некорректного ввода прошли. В отдельной серии из 100 запусков с одним муравьём и одной итерацией алгоритм в зависимости от начального состояния генератора либо находил допустимый маршрут, либо сообщал об исчерпании попыток.
